%-------------------------------------------------------------------------------
%	SECTION TITLE
%-------------------------------------------------------------------------------
\cvsection{Projects}


%-------------------------------------------------------------------------------
%	SUBSECTION TITLE
%-------------------------------------------------------------------------------
\begin{cventries}

%---------------------------------------------------------
     \cventry
    {Hackathon Winning Project}
    {ISO-Net : Multimodal Audio-Visual Target /speech Extraction }
    {}
    {}
    {
      \begin{cvitems}
        \item {Built a multimodal system to solve the cocktail-party problem, isolating a user-selected speaker's voice from noisy multi-speaker audio using a 4-microphone array.}
        \item {Combined spatial audio cues (GCC-PHAT) with face-conditioned visual embeddings inside a U-Net mask estimation network to enable face-selectable speech extraction.}
        \item {Achieved 9.31 dB SI-SDR (a 4.85 dB improvement over the raw mixture) on a hard -1 to 10 dB SNR test set, outperforming oracle delay-and-sum and MVDR beamformers.}
        \item {Won 1st place at KU Hackfeast 2025 and co-authored a resulting paper accepted for publication in JIEE.}
      \end{cvitems}
      ‎ 
    }
 %---------------------------------------------------------
     \cventry
    {Third year Project}
    {Human Presence Detection System Using Distributed Inference on Resource-Constrained Edge Devices}
    {}
    {}
    {
      \begin{cvitems}
        \item {Designed a distributed, tile-based computer vision system to enable real-time human presence detection on low-power microcontrollers.}
        \item {Implemented adaptive image tiling on an ESP32-CAM using Sobel-filter edge density to prioritize informative regions before transmission.}
        \item {Deployed a Local Binary CNN based model with 82\% accuracy on ESP32-S3 nodes, communicating over ESP-NOW for distributed inference.}
        \item {Aggregated distributed inference results on a master node to drive real-time detection status output via I2C LCD.}
        \item {Code : \href{https://github.com/Ishwor-git/hpd_distrributed_infrencing}{Github}}
      \end{cvitems}
      ‎ 
    }

    \cventry
    {AI Fellowship Graduation Project}
    {Detection of Hallucinations in Output of Reasoning Chains by LLMs }
    {}
    {}
    {
      \begin{cvitems}
        \item {Built lightweight classifiers to detect hallucinations in step-by-step reasoning chains generated by LLMs such as GPT.}
        \item {Benchmarked multiple architectures, from embedding-similarity ( 50 to 60\% accuracy  ) baselines to an LSTM model with an attention layer ( 70 to 80\% accuracy) , for hallucination detection.}
        \item {Curated custom dataset composed of various reasoning steps to solve language and logical problems that was used to train and evaluate}
        \item {code : \href{https://github.com/Ishwor-git/Detection-of-Hallucinations-in-Reasoning-Chains-by-LLMs}{Github}}
      \end{cvitems}
      ‎ 
    }
  \cventry
    {Ongoing Final Year project}
    {Spatiotemporal Graph Transformer for Clinically Interpretable Early Alzheimer’s Disease
Detection Using rs‑fMRI}
    {}
    {}
    {
      \begin{cvitems}
        \item {Designing a spatiotemporal graph transformer to classify Alzheimer’s Disease, Mild Cognitive Impairment, and Normal Controls from resting‑
state fMRI data in the ADNI dataset}
        \item {Combined Graph Attention Networks with a Temporal Transformer to jointly model spatial brain connectivity and its evolution over time.}
        \item {Applied self‑supervised contrastive pretraining (InfoNCE loss) to improve robustness on limited clinical neuroimaging samples.}
        \item {Building interpretability into the pipeline using attention analysis and Integrated Gradients, identifying the hippocampus and Default Mode Network
as key biomarkers.}
      \end{cvitems}
      ‎ 
    }
    

    


%---------------------------------------------------------
\end{cventries}
